Bổn Sâm có $n$ đồng xu \textbf{khác nhau về trọng lượng} và một cái cân hai bên.  
Mỗi khi anh ta đặt hai đồng xu $x$ và $y$ lên cân, trọng lượng tương đối giữa chúng sẽ được hiển thị --- tức là anh ta biết được giữa $x$ và $y$, đồng nào nặng hơn.

Gọi \textbf{cấp bậc} của đồng xu $x$ là số lượng đồng xu \textit{không nặng hơn} đồng xu $x$ (bao gồm cả chính nó).  
Do đó, đồng xu nhẹ nhất có cấp bậc là $1$, đồng nhẹ nhì là $2$, \ldots, và đồng nặng nhất có cấp bậc là $n$.

Sẽ có $m$ lần cân. Với mỗi đồng xu, hãy xác định sau \textbf{lần cân thứ $m$} thì có thể chắc chắn biết cấp bậc của đồng xu đó.
